% GENERATED FILE. Edit canonical lessons, then run npm run generate:books.
% canonical-source-hash: fnv1a64:36ac0d2864a61d24
% canonical-lessons: SA-C140-sucana, SA-C140-samayasarani, SA-C140-vijnapanam, SA-C140-deyakam, SA-C140-praptipatram, SA-R140-first-pass-the-future, SA-R140-second-pass-reading-notices

\chapter{A Notice, a Timetable, an Advert, a Bill, a Receipt}
\label{ch:sa-a-notice-a-timetable-an-advert-a-bill-a-receipt}
\hlchaptermodality{140}

\begin{chapteropening}
\textbf{By the end of this chapter:} \emph{I can read a notice, a timetable, an advert, a bill and a receipt, and act on what they say.}

Complete the last lesson of chapter 140: I can read a notice, a timetable, an advert, a bill and a receipt, and act on what they say.
\end{chapteropening}

\section[\texorpdfstring{\sk{सूचना} (sūcanā)}{सूचना (sūcanā)}]{\sk{सूचना} (sūcanā) — a notice}

\label{lesson:SA-C140-sucana}



\begin{quote}
Before the new one: say the Sanskrit for will be, then the Sanskrit for will go.
\end{quote}

\subsection*{You'll want to know: \sk{सूचना}}
\textbf{\sk{सूचना}} — \emph{sūcanā} — \textquotedblleft{}a notice\textquotedblright{}.

\emph{Read it:} the notice on the shop door — \textbf{\sk{अद्य} \sk{विश्रामः।}} — \textquotedblleft{}today, rest\textquotedblright{}

The shop is closed today, so come back \textbf{\sk{श्वः}}.

\begin{grammarlens}[title={What you've built}]
One more word for this chapter.
\end{grammarlens}

\subsection*{Your turn}
\begin{itemize}
\raggedright
  \item \emph{Say it:} \emph{sūcanā}
  \item \emph{Say it:} \emph{sūcanā}, once more
  \item \emph{Say it:} say \emph{sūcanā}
  \item \emph{Recall:} say the Sanskrit for will be, then the Sanskrit for will go, then say \emph{sūcanā} again
\end{itemize}

\begin{tcolorbox}[breakable,colback=teal!4,colframe=teal!35!black,title={Before you move on}]
What does \sk{सूचना} mean? (\textquotedblleft{}A notice\textquotedblright{}.) Say it once more.
\end{tcolorbox}
\section[\texorpdfstring{\sk{समयसारणी} (samayasāraṇī)}{समयसारणी (samayasāraṇī)}]{\sk{समयसारणी} (samayasāraṇī) — a timetable}

\label{lesson:SA-C140-samayasarani}



\begin{quote}
Before the new one: say the Sanskrit for will go, then the Sanskrit for a notice.
\end{quote}

\subsection*{You'll want to know: \sk{समयसारणी}}
\textbf{\sk{समयसारणी}} — \emph{samayasāraṇī} — \textquotedblleft{}a timetable\textquotedblright{}.

\emph{Read it:} the timetable at the station — \textbf{\sk{धूमयानम्} — \sk{प्रातः} \sk{दश।}}

The train leaves at ten in the morning: be at the \textbf{\sk{स्थानकम्}} before ten.

\begin{grammarlens}[title={What you've built}]
One more word for this chapter.
\end{grammarlens}

\subsection*{Your turn}
\begin{itemize}
\raggedright
  \item \emph{Say it:} \emph{samayasāraṇī}
  \item \emph{Say it:} \emph{samayasāraṇī}, once more
  \item \emph{Say it:} say \emph{samayasāraṇī}
  \item \emph{Recall:} say the Sanskrit for will go, then the Sanskrit for a notice, then say \emph{samayasāraṇī} again
\end{itemize}

\begin{tcolorbox}[breakable,colback=teal!4,colframe=teal!35!black,title={Before you move on}]
What does \sk{समयसारणी} mean? (\textquotedblleft{}A timetable\textquotedblright{}.) Say it once more.
\end{tcolorbox}
\section[\texorpdfstring{\sk{विज्ञापनम्} (vijñāpanam)}{विज्ञापनम् (vijñāpanam)}]{\sk{विज्ञापनम्} (vijñāpanam) — an advert}

\label{lesson:SA-C140-vijnapanam}



\begin{quote}
Before the new one: say the Sanskrit for a notice, then the Sanskrit for a timetable.
\end{quote}

\subsection*{You'll want to know: \sk{विज्ञापनम्}}
\textbf{\sk{विज्ञापनम्}} — \emph{vijñāpanam} — \textquotedblleft{}an advert\textquotedblright{}.

\emph{Read it:} the advert — \textbf{\sk{प्रकोष्ठः} \sk{रिक्तः।} \sk{मूल्यम्} — \sk{मासे} \sk{सहस्रम्।}}

A room is empty, at a thousand a month.

\begin{grammarlens}[title={What you've built}]
One more word for this chapter.
\end{grammarlens}

\subsection*{Your turn}
\begin{itemize}
\raggedright
  \item \emph{Say it:} \emph{vijñāpanam}
  \item \emph{Say it:} \emph{vijñāpanam}, once more
  \item \emph{Say it:} say \emph{vijñāpanam}
  \item \emph{Recall:} say the Sanskrit for a notice, then the Sanskrit for a timetable, then say \emph{vijñāpanam} again
\end{itemize}

\begin{tcolorbox}[breakable,colback=teal!4,colframe=teal!35!black,title={Before you move on}]
What does \sk{विज्ञापनम्} mean? (\textquotedblleft{}An advert\textquotedblright{}.) Say it once more.
\end{tcolorbox}
\section[\texorpdfstring{\sk{देयकम्} (deyakam)}{देयकम् (deyakam)}]{\sk{देयकम्} (deyakam) — a bill (to pay)}

\label{lesson:SA-C140-deyakam}



\begin{quote}
Before the new one: say the Sanskrit for a timetable, then the Sanskrit for an advert.
\end{quote}

\subsection*{You'll want to know: \sk{देयकम्}}
\textbf{\sk{देयकम्}} — \emph{deyakam} — \textquotedblleft{}a bill (to pay)\textquotedblright{}.

\emph{Read it:} the bill at the restaurant — \textbf{\sk{भोजनालयः} — \sk{देयकम्} — \sk{शतम्।}}

A hundred to pay before you go.

\begin{grammarlens}[title={What you've built}]
One more word for this chapter.
\end{grammarlens}

\subsection*{Your turn}
\begin{itemize}
\raggedright
  \item \emph{Say it:} \emph{deyakam}
  \item \emph{Say it:} \emph{deyakam}, once more
  \item \emph{Say it:} say \emph{deyakam}
  \item \emph{Recall:} say the Sanskrit for a timetable, then the Sanskrit for an advert, then say \emph{deyakam} again
\end{itemize}

\begin{tcolorbox}[breakable,colback=teal!4,colframe=teal!35!black,title={Before you move on}]
What does \sk{देयकम्} mean? (\textquotedblleft{}A bill (to pay)\textquotedblright{}.) Say it once more.
\end{tcolorbox}
\section[\texorpdfstring{\sk{प्राप्तिपत्रम्} (prāptipatram)}{प्राप्तिपत्रम् (prāptipatram)}]{\sk{प्राप्तिपत्रम्} (prāptipatram) — a receipt}

\label{lesson:SA-C140-praptipatram}



\begin{quote}
Before the new one: say the Sanskrit for an advert, then the Sanskrit for a bill.
\end{quote}

\subsection*{You'll want to know: \sk{प्राप्तिपत्रम्}}
\textbf{\sk{प्राप्तिपत्रम्}} — \emph{prāptipatram} — \textquotedblleft{}a receipt\textquotedblright{}.

\emph{Read it:} the receipt — \textbf{\sk{प्रकोष्ठः} — \sk{मूल्यम्} — \sk{सहस्रम्।}}

The thousand for the room is paid: keep it.

\begin{grammarlens}[title={What you've built}]
One more word for this chapter.
\end{grammarlens}

\subsection*{Your turn}
\begin{itemize}
\raggedright
  \item \emph{Say it:} \emph{prāptipatram}
  \item \emph{Say it:} \emph{prāptipatram}, once more
  \item \emph{Say it:} say \emph{prāptipatram}
  \item \emph{Recall:} say the Sanskrit for an advert, then the Sanskrit for a bill, then say \emph{prāptipatram} again
\end{itemize}

\begin{tcolorbox}[breakable,colback=teal!4,colframe=teal!35!black,title={Before you move on}]
What does \sk{प्राप्तिपत्रम्} mean? (\textquotedblleft{}A receipt\textquotedblright{}.) Say it once more.
\end{tcolorbox}
\section[(dialogue)]{The future — a first pass}

\label{lesson:SA-R140-first-pass-the-future}



\begin{quote}
Say the words for a bill and a receipt.
\end{quote}

\subsection*{Your turn}
\begin{itemize}
\raggedright
  \item \emph{Say it:} the future, tomorrow, coming, will be, will go
\end{itemize}

\begin{tcolorbox}[breakable,colback=teal!4,colframe=teal!35!black,title={Before you move on}]
The future? (\textbf{\sk{भविष्यत्}}, \emph{bhaviṣyat}.) Tomorrow? (\textbf{\sk{श्वः}}, \emph{śvaḥ}.) Coming? (\textbf{\sk{आगामी}}, \emph{āgāmī}.) Will be? (\textbf{\sk{भविष्यति}}, \emph{bhaviṣyati}.) Will go? (\textbf{\sk{गमिष्यति}}, \emph{gamiṣyati}.)
\end{tcolorbox}
\section[(dialogue)]{Reading notices — a second pass}

\label{lesson:SA-R140-second-pass-reading-notices}



\begin{quote}
Say the words for a notice and a timetable.
\end{quote}

\subsection*{Your turn}
\begin{itemize}
\raggedright
  \item \emph{Say it:} a notice, a timetable, an advert, a bill, a receipt
\end{itemize}

\begin{tcolorbox}[breakable,colback=teal!4,colframe=teal!35!black,title={Before you move on}]
A notice? (\textbf{\sk{सूचना}}, \emph{sūcanā}.) A timetable? (\textbf{\sk{समयसारणी}}, \emph{samayasāraṇī}.) An advert? (\textbf{\sk{विज्ञापनम्}}, \emph{vijñāpanam}.) A bill? (\textbf{\sk{देयकम्}}, \emph{deyakam}.) A receipt? (\textbf{\sk{प्राप्तिपत्रम्}}, \emph{prāptipatram}.)
\end{tcolorbox}
